\documentclass[aps,preprint]{revtex4}%
\usepackage{amssymb}
\usepackage{amsfonts}
\usepackage{amsmath}
\usepackage{graphicx}%
\setcounter{MaxMatrixCols}{30}
%TCIDATA{OutputFilter=latex2.dll}
%TCIDATA{Version=5.50.0.2960}
%TCIDATA{CSTFile=revtex4.cst}
%TCIDATA{Created=Thursday, June 27, 2013 17:39:45}
%TCIDATA{LastRevised=Friday, October 17, 2014 09:47:09}
%TCIDATA{<META NAME="GraphicsSave" CONTENT="32">}
%TCIDATA{<META NAME="SaveForMode" CONTENT="1">}
%TCIDATA{BibliographyScheme=Manual}
%TCIDATA{<META NAME="DocumentShell" CONTENT="Articles\SW\REVTeX 4">}
%BeginMSIPreambleData
\providecommand{\U}[1]{\protect\rule{.1in}{.1in}}
%EndMSIPreambleData
\newtheorem{theorem}{Theorem}
\newtheorem{acknowledgement}[theorem]{Acknowledgement}
\newtheorem{algorithm}[theorem]{Algorithm}
\newtheorem{axiom}[theorem]{Axiom}
\newtheorem{claim}[theorem]{Claim}
\newtheorem{conclusion}[theorem]{Conclusion}
\newtheorem{condition}[theorem]{Condition}
\newtheorem{conjecture}[theorem]{Conjecture}
\newtheorem{corollary}[theorem]{Corollary}
\newtheorem{criterion}[theorem]{Criterion}
\newtheorem{definition}[theorem]{Definition}
\newtheorem{example}[theorem]{Example}
\newtheorem{exercise}[theorem]{Exercise}
\newtheorem{lemma}[theorem]{Lemma}
\newtheorem{notation}[theorem]{Notation}
\newtheorem{problem}[theorem]{Problem}
\newtheorem{proposition}[theorem]{Proposition}
\newtheorem{remark}[theorem]{Remark}
\newtheorem{solution}[theorem]{Solution}
\newtheorem{summary}[theorem]{Summary}
\newenvironment{proof}[1][Proof]{\noindent\textbf{#1.} }{\ \rule{0.5em}{0.5em}}
\begin{document}
\preprint{HEP/123-qed}
\title[Short title for running header]{REV\TeX4}
\author{First Author}
\affiliation{My Institution}
\author{Second Author}
\affiliation{My Institution}
\author{Third Author}
\affiliation{Other Institution}
\keywords{one two three}
\pacs{PACS number}

\begin{abstract}
Shell document for REV\TeX{} 4.

\end{abstract}
\volumeyear{year}
\volumenumber{number}
\issuenumber{number}
\eid{identifier}
\date[Date text]{date}
\received[Received text]{date}

\revised[Revised text]{date}

\accepted[Accepted text]{date}

\published[Published text]{date}

\startpage{101}
\endpage{102}
\maketitle
\tableofcontents


\subsection{Representations}

The Pauli matrices are
\begin{align}
\mathbb{\sigma}_{0}  &  =\left(
\begin{array}
[c]{cc}%
1 & 0\\
0 & 1
\end{array}
\right)  =\mathbb{I}_{2}\\
\mathbb{\sigma}_{1}  &  =\left(
\begin{array}
[c]{cc}%
0 & 1\\
1 & 0
\end{array}
\right)  =\mathbb{\sigma}_{x}\\
\mathbb{\sigma}_{2}  &  =\left(
\begin{array}
[c]{cc}%
0 & -i\\
i & 0
\end{array}
\right)  =\mathbb{\sigma}_{y}\\
\mathbb{\sigma}_{3}  &  =\left(
\begin{array}
[c]{cc}%
1 & 0\\
0 & -1
\end{array}
\right)  =\mathbb{\sigma}_{z}%
\end{align}
Representations of the complex Clifford algebras in terms of the Pauli
matrices are given by%
\begin{align}
\operatorname*{Cl}\left(  2m,\mathbb{C}\right)   &  \simeq M\left(
2^{m},\mathbb{C}\right) \\
\operatorname*{Cl}\left(  2m-1,\mathbb{C}\right)   &  \simeq M\left(
2^{m-1},\mathbb{C}\right)  \oplus M\left(  2^{m-1},\mathbb{C}\right) \\
M\left(  2^{m},\mathbb{C}\right)   &  :\mathbb{C}^{2m}\rightarrow
\mathbb{C}^{2m}\\
\mathbb{C}^{2m}  &  \simeq\,\overset{\text{m times}}{\overbrace{\,\,\mathbb{C}%
^{2}\otimes\cdots\otimes\mathbb{C}^{2}}}\\
\dim\left\{  \operatorname*{Cl}\left(  2m,\mathbb{C}\right)  \right\}   &
=2^{2m}%
\end{align}


and we can defines annihilation and creation operators belonging to $%
%TCIMACRO{\dbigotimes \limits^{m}}%
%BeginExpansion
{\displaystyle\bigotimes\limits^{m}}
%EndExpansion
M\left(  2,\mathbb{C}\right)  $ as
\begin{align}
a_{k}  &  =\,\overset{m-k-1}{\overbrace{\mathbb{I}_{2}\otimes\cdots
\otimes\mathbb{I}_{2}}}\otimes a\otimes\overset{k}{\overbrace{\mathbb{\sigma
}_{z}\otimes\cdots\otimes\mathbb{\sigma}_{z}}}\nonumber\\
a_{k}^{\dagger}  &  =\overset{m-k-1}{\overbrace{\,\mathbb{I}_{2}\otimes
\cdots\otimes\mathbb{I}_{2}}}\otimes a^{\dagger}\otimes
\overset{k}{\overbrace{\mathbb{\sigma}_{z}\otimes\cdots\otimes\mathbb{\sigma
}_{z}}}\nonumber\\
1  &  \leq k\leq m,
\end{align}
where
\begin{equation}
a=\left(
\begin{array}
[c]{cc}%
0 & 1\\
0 & 0
\end{array}
\right)
\end{equation}
Hence
\begin{align}
x_{k}  &  =\frac{1}{\sqrt{2}}\left(  a_{k}^{\dagger}+a_{k}\right)  =\frac
{1}{\sqrt{2}}\left(  \overset{m-k-1}{\overbrace{\,\mathbb{I}_{2}\otimes
\cdots\otimes\mathbb{I}_{2}}}\otimes a^{\dagger}\otimes
\overset{k}{\overbrace{\mathbb{\sigma}_{z}\otimes\cdots\otimes\mathbb{\sigma
}_{z}}}+\,\overset{m-k-1}{\overbrace{\mathbb{I}_{2}\otimes\cdots
\otimes\mathbb{I}_{2}}}\otimes a\otimes\overset{k}{\overbrace{\mathbb{\sigma
}_{z}\otimes\cdots\otimes\mathbb{\sigma}_{z}}}\right) \\
&  =\frac{1}{\sqrt{2}}\left(  \overset{m-k-1}{\overbrace{\,\mathbb{I}%
_{2}\otimes\cdots\otimes\mathbb{I}_{2}}}\otimes\left(  a^{\dagger}+a\right)
\otimes\overset{k}{\overbrace{\mathbb{\sigma}_{z}\otimes\cdots\otimes
\mathbb{\sigma}_{z}}}\right) \\
&  =\frac{1}{\sqrt{2}}\left(  \overset{m-k-1}{\overbrace{\,\mathbb{I}%
_{2}\otimes\cdots\otimes\mathbb{I}_{2}}}\otimes\left(
\begin{array}
[c]{cc}%
0 & 1\\
1 & 0
\end{array}
\right)  \otimes\overset{k}{\overbrace{\mathbb{\sigma}_{z}\otimes\cdots
\otimes\mathbb{\sigma}_{z}}}\right) \\
x_{k+r}  &  =\frac{i}{\sqrt{2}}\left(  a_{k}^{\dagger}-ia_{k}\right)
=\frac{i}{\sqrt{2}}\left(  \overset{m-k-1}{\overbrace{\,\mathbb{I}_{2}%
\otimes\cdots\otimes\mathbb{I}_{2}}}\otimes\left(  a^{\dagger}-ia\right)
\otimes\overset{k}{\overbrace{\mathbb{\sigma}_{z}\otimes\cdots\otimes
\mathbb{\sigma}_{z}}}\right) \\
&  =\frac{i}{\sqrt{2}}\left(  \overset{m-k-1}{\overbrace{\,\mathbb{I}%
_{2}\otimes\cdots\otimes\mathbb{I}_{2}}}\otimes\left(
\begin{array}
[c]{cc}%
0 & -i\\
1 & 0
\end{array}
\right)  \otimes\overset{k}{\overbrace{\mathbb{\sigma}_{z}\otimes\cdots
\otimes\mathbb{\sigma}_{z}}}\right)
\end{align}


\subsubsection{Examples}

\paragraph{m=1 then}

\subparagraph{ even}

These are the Pauli matrices themselves i.e.%
\begin{align}
\operatorname*{Cl}\left(  2,\mathbb{C}\right)   &  \simeq M\left(
2,\mathbb{C}\right) \\
M\left(  2,\mathbb{C}\right)   &  :\mathbb{C}^{2}\rightarrow\mathbb{C}^{2}%
\end{align}


\subparagraph{odd%
\begin{equation}
\operatorname*{Cl}\left(  1,\mathbb{C}\right)  \simeq M\left(  2^{0}%
,\mathbb{C}\right)  \oplus M\left(  2^{0},\mathbb{C}\right)  \simeq
\mathbb{C}\oplus\mathbb{C},
\end{equation}
}

\paragraph{m=2 then}

\subparagraph{even}%

\begin{align}
\operatorname*{Cl}\left(  4,\mathbb{C}\right)   &  \simeq M\left(
4,\mathbb{C}\right) \\
M\left(  4,\mathbb{C}\right)   &  :\mathbb{C}^{4}\rightarrow\mathbb{C}^{4}%
\end{align}
The creator and annihilator are given by
\begin{align}
a_{1}  &  =\,\overset{2-1-1}{\overbrace{\mathbb{I}_{2}\otimes\cdots
\otimes\mathbb{I}_{2}}}\otimes a\otimes\overset{1}{\overbrace{\mathbb{\sigma
}_{z}\otimes\cdots\otimes\mathbb{\sigma}_{z}}}\,=a\otimes\mathbb{\sigma}_{z}\\
a_{1}^{\dagger}  &  =\overset{2-1-1}{\overbrace{\,\mathbb{I}_{2}\otimes
\cdots\otimes\mathbb{I}_{2}}}\otimes a^{\dagger}\otimes
\overset{1}{\overbrace{\mathbb{\sigma}_{z}\otimes\cdots\otimes\mathbb{\sigma
}_{z}}}\,=a^{\dagger}\otimes\mathbb{\sigma}_{z}.
\end{align}
The dimension of this Fock space is
\begin{equation}
\dim\left\{  \mathcal{H}_{\mathcal{F}}\right\}  =4=1+2+1=2^{2}=2^{2r}%
\end{equation}
and of this Fock algebra is
\begin{equation}
\dim\left\{  \mathcal{F}\right\}  =16=\left(  2^{2}\right)  ^{2}=\left(
2^{2r}\right)  ^{2}=2^{4}=2^{4r}%
\end{equation}
as the number of degrees of freedom $\dim\left\{  \mathcal{H}^{1}\right\}
=r=1.$The scalar product is of rank $=2$. The dimension of the Clifford
algebra is thus $2$ and thus represented by matrices of order $2^{2}=4$

\subparagraph{odd%
\begin{equation}
\operatorname*{Cl}\left(  3,\mathbb{C}\right)  \simeq M\left(  2,\mathbb{C}%
\right)  \oplus M\left(  2,\mathbb{C}\right)
\end{equation}
}

\subsection{Basis}

A basis for a Fock representation of the Clifford Algebra $\operatorname*{Cl}%
\left(  4r,\mathbb{C}\right)  $ $\simeq M\left(  2^{2r},\mathbb{C}\right)  $
is
\begin{equation}
I_{\mathcal{F}},\left\{  x_{j_{1}},x_{j_{1}}x_{j_{2}},\cdots,x_{j_{1}}\cdots
x_{j_{2r}}\right\}  ;1\leq j_{1}<\cdots<j_{2r}\leq2r
\end{equation}
Number of basis elements and thus dimension of $\mathcal{F}$ is
\begin{equation}
2^{4r}=\sum_{0\leq j\leq2r}\binom{2r}{j}.
\end{equation}


\subsection{Trace}

The multiplicative unit is $I_{\mathcal{F}}$ and has the property
\begin{equation}
\left\langle I_{\mathcal{F}}\right\rangle _{0}=I_{\mathcal{F}}%
\end{equation}
and
\begin{equation}
\left\langle \alpha I_{\mathcal{F}}\right\rangle _{0}=\alpha I_{\mathcal{F}}%
\end{equation}
One can introduce a linear functional, $\operatorname*{tr},$%
\begin{align}
\operatorname*{tr} &  :\mathcal{F\rightarrow}\mathbb{C}\\
\operatorname*{tr}\left\{  A\right\}   &  \equiv\left\langle \left\langle
A\right\rangle \right\rangle _{0}%
\end{align}
that has the properties
\begin{align}
\left\langle A\right\rangle _{0} &  =\operatorname*{tr}\left\{  A\right\}
I_{\mathcal{F}}\\
\operatorname*{tr}\left\{  I_{\mathcal{F}}\right\}   &  =1\\
\operatorname*{tr}\left\{  \alpha I_{\mathcal{F}}\right\}   &  =\alpha\\
\operatorname*{tr}\left\{  A+B\right\}   &  =\operatorname*{tr}\left\{
A\right\}  +\operatorname*{tr}\left\{  B\right\}  \\
\operatorname*{tr}\left\{  AB\right\}   &  =\operatorname*{tr}\left\{
BA\right\}  \\
\operatorname*{tr}\left\{  ABC\right\}   &  =\operatorname*{tr}\left\{
CAB\right\}  =\operatorname*{tr}\left\{  BCA\right\}  \\
\left\langle \left\langle A^{\dagger}B\right\rangle \right\rangle _{0} &
=\operatorname*{tr}\left\{  A^{\dagger}B\right\}  I_{\mathcal{F}}\\
\left\Vert A\right\Vert ^{2} &  =\left\langle \left\langle A^{\dagger
}A\right\rangle \right\rangle _{0}=\operatorname*{tr}\left\{  A^{\dagger
}A\right\}  \\
\left\vert A\right\vert ^{2} &  =\left\langle \left\langle A^{t}A\right\rangle
\right\rangle _{0}=\operatorname*{tr}\left\{  A^{t}A\right\}
\end{align}
and
\begin{equation}
\operatorname*{tr}\left\{  A\right\}  =2^{-2r}\operatorname*{Tr}\left\{
A\right\}  ,
\end{equation}
where $\operatorname*{Tr}$ is the normal trace on the Fock representation.

\subsection{Volume Element}%

\begin{equation}
\mathbf{I}=x_{1}\wedge\cdots\wedge x_{2r}%
\end{equation}
If $r=\infty$ then this definition needs careful analysis.

\subsection{Geometric Product}

The geometric product between two general elements is defined as
\begin{equation}
AB=\sum_{0\leq m,n\leq2r}A_{m}B_{n},
\end{equation}
where
\begin{equation}
A_{m}B_{n}=\sum_{j=0}^{\min\left(  m,n\right)  }\left\langle A_{m}%
B_{n}\right\rangle _{\left\vert m-n\right\vert +2j}.
\end{equation}
E.g. let
\begin{align}
A  &  =A_{0}+A_{1}+A_{2}\\
B  &  =B_{0}+B_{1}+B_{2}+B_{3}+B_{4}%
\end{align}
then
\begin{equation}
AB=A_{0}\left(  B_{0}+B_{1}+B_{2}+B_{3}+B_{4}\right)  +A_{1}\left(
B_{0}+B_{1}+B_{2}+B_{3}+B_{4}\right)  +A_{2}\left(  B_{0}+B_{1}+B_{2}%
+B_{3}+B_{4}\right)  .
\end{equation}
$\lceil$%
\begin{equation}
AB=\sum_{1\leq j,k,l,m,n,o\leq2r}A_{jk}B_{lmno}x_{j}x_{k}x_{l}x_{m}x_{n}x_{o}%
\end{equation}
The parts containing $A_{1}$ are given by
\begin{align}
A_{1}B_{1}  &  =\sum_{j=0}^{1}\left\langle A_{1}B_{1}\right\rangle
_{2j}=\left\langle A_{1}B_{1}\right\rangle _{0}+\left\langle A_{1}%
B_{1}\right\rangle _{2}\\
A_{1}B_{2}  &  =\sum_{j=0}^{1}\left\langle A_{1}B_{2}\right\rangle
_{1+2j}=\left\langle A_{1}B_{2}\right\rangle _{1}+\left\langle A_{1}%
B_{2}\right\rangle _{3}\\
A_{1}B_{3}  &  =\sum_{j=0}^{1}\left\langle A_{1}B_{3}\right\rangle
_{2+2j}=\left\langle A_{1}B_{3}\right\rangle _{2}+\left\langle A_{1}%
B_{3}\right\rangle _{4}\\
A_{1}B_{4}  &  =\sum_{j=0}^{1}\left\langle A_{1}B_{4}\right\rangle
_{3+2j}=\left\langle A_{1}B_{4}\right\rangle _{3}+\left\langle A_{1}%
B_{4}\right\rangle _{5}%
\end{align}
Expanding the operators in a basis gives
\begin{align}
\left\langle A_{1}B_{1}\right\rangle _{0}  &  =\left\langle \sum_{1\leq
j,k\leq2r}A_{1j}B_{1k}x_{j}x_{k}\right\rangle _{0}=\sum_{1\leq j\leq2r}%
A_{1j}B_{1k}x_{j}^{2}\\
\left\langle A_{1}B_{1}\right\rangle _{2}  &  =\left\langle \sum_{1\leq
j,k\leq2r}A_{1j}B_{1k}x_{j}x_{k}\right\rangle _{2}=\sum_{1\leq j<k\leq
2r}\left(  A_{1j}B_{1k}-A_{1k}B_{1j}\right)  x_{j}x_{k}=\sum_{1\leq j<k\leq
2r}\left\vert
\begin{array}
[c]{cc}%
A_{1j} & A_{1k}\\
B_{1j} & B_{1k}%
\end{array}
\right\vert x_{j}x_{k}%
\end{align}%
\begin{align}
\left\langle A_{1}B_{2}\right\rangle _{1}  &  =\left\langle \frac{1}{2}%
\sum_{1\leq j,k\neq l\leq2r}A_{1j}B_{2kl}x_{j}x_{k}x_{l}\right\rangle
_{1}=\frac{1}{2}\sum_{1\leq j,k\leq2r}\left(  A_{1j}B_{2jk}-A_{1j}%
B_{2kj}\right)  x_{j}^{2}x_{k}\\
\left\langle A_{1}B_{2}\right\rangle _{3}  &  =\left\langle \frac{1}{2}%
\sum_{1\leq j,k\neq l\leq2r}A_{1j}B_{2kl}x_{j}x_{k}x_{l}\right\rangle
_{3}=\sum_{1\leq j<k<l\leq2r}\left(  A_{1j}B_{2kl}-A_{1k}B_{2jl}+A_{1l}%
B_{kj}\right)  x_{j}x_{k}x_{l}%
\end{align}%
\begin{align}
\left\langle A_{1}B_{3}\right\rangle _{2}  &  =\left\langle \frac{1}{6}%
\sum_{1\leq j,k\neq l\neq m\leq2r}A_{1j}B_{3klm}x_{j}x_{k}x_{l}x_{m}%
\right\rangle _{2}=\frac{1}{3}\left\{  \sum_{1\leq j<k,l\leq2r}\left(
A_{1j}B_{3kll}-A_{1j}B_{3lkl}+A_{1l}B_{3ljk}\right)  x_{j}x_{k}x_{l}%
^{2}\right\} \\
\left\langle A_{1}B_{3}\right\rangle _{4}  &  =\left\langle \frac{1}{6}%
\sum_{1\leq j,k\neq l\neq m\leq2r}A_{1j}B_{3klm}x_{j}x_{k}x_{l}x_{m}%
\right\rangle _{4}=\sum_{1\leq j<k<l<m\leq2r}\left(  A_{1j}B_{3klm}%
-A_{1k}B_{3jlm}+A_{1l}B_{3kjm}-A_{1m}B_{3klj}\right)  x_{j}x_{k}x_{l}x_{m}%
\end{align}%
\begin{align}
\left\langle A_{1}B_{4}\right\rangle _{3}  &  =\left\langle \frac{1}{24}%
\sum_{1\leq j_{1},j_{2}\neq j_{3}\neq j_{4}\neq j_{5}\leq2r}A_{1j_{1}%
}B_{4j_{2}j_{3}j_{4}j_{5}}x_{j_{1}}x_{j_{2}}x_{j_{3}}x_{j_{4}}x_{j_{5}%
}\right\rangle _{3}=\\
&  \frac{1}{24}\sum_{1\leq j_{1},j_{2}\neq j_{3}\neq j_{4}\neq j_{5}\leq
2r}A_{1j_{1}}B_{4j_{2}j_{3}j_{4}j_{5}}\left\{  \delta_{j_{1}j_{2}}x_{j_{3}%
}\wedge x_{j_{4}}\wedge x_{j_{5}}-\delta_{j_{1}j_{3}}x_{j_{2}}\wedge x_{j_{4}%
}\wedge x_{j_{5}}+\delta_{j_{1}j_{4}}x_{j_{2}}\wedge x_{j_{3}}\wedge x_{j_{5}%
}+\delta_{j_{2}j_{3}}x_{j_{1}}\wedge x_{j_{4}}\wedge x_{j_{5}}\right. \\
&  -\delta_{j_{2}j_{4}}x_{j_{1}}\wedge x_{j_{3}}\wedge x_{j_{5}}+\delta
_{j_{3}j_{4}}x_{j_{1}}\wedge x_{j_{2}}\wedge x_{j_{5}}+\delta_{j_{4}j_{5}%
}x_{j_{1}}\wedge x_{j_{2}}\wedge x_{j_{3}}-\delta_{j_{3}j_{5}}x_{j_{1}}\wedge
x_{j_{2}}\wedge x_{j_{4}}\\
&  +\left.  \delta_{j_{2}j_{5}}x_{j_{1}}\wedge x_{j_{3}}\wedge x_{j_{4}%
}-\delta_{j_{1}j_{5}}x_{j_{2}}\wedge x_{j_{3}}\wedge x_{j_{4}}\right\} \\
\left\langle A_{1}B_{4}\right\rangle _{5}  &  =\left\langle \frac{1}{24}%
\sum_{1\leq j_{1}\neq j_{2}\neq j_{3}\neq j_{4}\neq j_{5}\leq2r}A_{1j_{1}%
}B_{4j_{2}j_{3}j_{4}j_{5}}x_{j_{1}}x_{j_{2}}x_{j_{3}}x_{j_{4}}x_{j_{5}%
}\right\rangle _{5}=\frac{1}{24}\sum_{1\leq j_{1}\neq j_{2}\neq j_{3}\neq
j_{4}\neq j_{5}\leq2r}A_{1j_{1}}B_{4j_{2}j_{3}j_{4}j_{5}}x_{j_{1}}\wedge
x_{j_{2}}\wedge x_{j_{3}}\wedge x_{j_{4}}\wedge x_{j_{5}}%
\end{align}%
\begin{align}
&  x_{j_{1}}x_{j_{2}}x_{j_{3}}x_{j_{4}}x_{j_{5}}=x_{j_{1}}\wedge x_{j_{2}%
}\wedge x_{j_{3}}\wedge x_{j_{4}}\wedge x_{j_{5}}\nonumber\\
&  +\Delta_{j_{1}j_{2}}x_{j_{3}}\wedge x_{j_{4}}\wedge x_{j_{5}}-\Delta
_{j_{1}j_{3}}x_{j_{2}}\wedge x_{j_{4}}\wedge x_{j_{5}}+\Delta_{j_{1}j_{4}%
}x_{j_{2}}\wedge x_{j_{3}}\wedge x_{j_{5}}+\Delta_{j_{2}j_{3}}x_{j_{1}}\wedge
x_{j_{4}}\wedge x_{j_{5}}\nonumber\\
&  -\Delta_{j_{2}j_{4}}x_{j_{1}}\wedge x_{j_{3}}\wedge x_{j_{5}}+\Delta
_{j_{3}j_{4}}x_{j_{1}}\wedge x_{j_{2}}\wedge x_{j_{5}}+\Delta_{j_{4}j_{5}%
}x_{j_{1}}\wedge x_{j_{2}}\wedge x_{j_{3}}-\Delta_{j_{3}j_{5}}x_{j_{1}}\wedge
x_{j_{2}}\wedge x_{j_{4}}\nonumber\\
&  +\Delta_{j_{2}j_{5}}x_{j_{1}}\wedge x_{j_{3}}\wedge x_{j_{4}}-\Delta
_{j_{1}j_{5}}x_{j_{2}}\wedge x_{j_{3}}\wedge x_{j_{4}}\nonumber\\
&  +\operatorname*{Pf}\left(  \mathbf{\Delta}\right)  _{j_{2}j_{3}j_{4}j_{5}%
}x_{j_{1}}-\operatorname*{Pf}\left(  \mathbf{\Delta}\right)  _{j_{1}j_{3}%
j_{4}j_{5}}x_{j}+\operatorname*{Pf}\left(  \mathbf{\Delta}\right)
_{j_{1}j_{2}j_{4}j_{5}}x_{j_{3}}-\operatorname*{Pf}\left(  \mathbf{\Delta
}\right)  _{j_{1}j_{2}j_{3}j_{5}}x_{j_{4}}+\operatorname*{Pf}\left(
\mathbf{\Delta}\right)  _{j_{1}j_{2}j_{3}j_{4}}x_{j_{5}},
\end{align}%
\begin{equation}
\operatorname*{Pf}\left(  \mathbf{\Delta}\right)  _{j_{1}j_{2}j_{3}j_{4}%
}=\Delta_{j_{1}j_{2}}\Delta_{j_{3}j_{4}}-\Delta_{j_{1}j_{3}}\Delta_{j_{2}%
j_{4}}+\Delta_{j_{1}j_{4}}\Delta_{j_{2}j_{3}}.
\end{equation}
The parts containing $A_{2}$ are given by
\begin{align}
A_{2}B_{1}  &  =\sum_{j=0}^{1}\left\langle A_{2}B_{1}\right\rangle
_{1+2j}=\left\langle A_{2}B_{1}\right\rangle _{1}+\left\langle A_{2}%
B_{1}\right\rangle _{3}\\
A_{2}B_{2}  &  =\sum_{j=0}^{2}\left\langle A_{2}B_{2}\right\rangle
_{2j}=\left\langle A_{2}B_{2}\right\rangle _{0}+\left\langle A_{2}%
B_{2}\right\rangle _{2}+\left\langle A_{2}B_{2}\right\rangle _{4}\\
A_{2}B_{3}  &  =\sum_{j=0}^{2}\left\langle A_{2}B_{3}\right\rangle
_{1+2j}=\left\langle A_{2}B_{3}\right\rangle _{1}+\left\langle A_{2}%
B_{3}\right\rangle _{3}+\left\langle A_{2}B_{3}\right\rangle _{5}\\
A_{2}B_{4}  &  =\sum_{j=0}^{2}\left\langle A_{2}B_{4}\right\rangle
_{2+2j}=\left\langle A_{2}B_{4}\right\rangle _{2}+\left\langle A_{2}%
B_{4}\right\rangle _{4}+\left\langle A_{2}B_{4}\right\rangle _{6},
\end{align}
where
\begin{align}
\left\langle A_{2}B_{1}\right\rangle _{1}  &  =\frac{1}{2}\left\langle
\sum_{1\leq j\neq k,l\leq2r}A_{2jk}B_{1l}x_{j}x_{k}x_{l}\right\rangle _{1}\\
\left\langle A_{2}B_{1}\right\rangle _{3}  &  =\frac{1}{2}\left\langle
\sum_{1\leq j\neq k,l\leq2r}A_{2jk}B_{1l}x_{j}x_{k}x_{l}\right\rangle _{3}%
\end{align}
%

\begin{align}
\left\langle A_{2}B_{2}\right\rangle _{0}  &  =\frac{1}{4}\left\langle
\sum_{1\leq j\neq k,l\neq m\leq2r}A_{2jk}B_{2lm}x_{j}x_{k}x_{l}x_{m}%
\right\rangle _{0}\\
\left\langle A_{2}B_{2}\right\rangle _{2}  &  =\frac{1}{4}\left\langle
\sum_{1\leq j\neq k,l\neq m\leq2r}A_{2jk}B_{2lm}x_{j}x_{k}x_{l}x_{m}%
\right\rangle _{2}\\
\left\langle A_{2}B_{2}\right\rangle _{4}  &  =\frac{1}{4}\left\langle
\sum_{1\leq j\neq k,l\neq m\leq2r}A_{2jk}B_{2lm}x_{j}x_{k}x_{l}x_{m}%
\right\rangle _{4}%
\end{align}
%

\begin{align}
\left\langle A_{2}B_{3}\right\rangle _{1}  &  =\frac{1}{12}\left\langle
\sum_{1\leq j\neq k,l\neq m\neq n\leq2r}A_{2jk}B_{3lmn}x_{j}x_{k}x_{l}%
x_{m}x_{n}\right\rangle _{1}\\
\left\langle A_{2}B_{3}\right\rangle _{3}  &  =\frac{1}{12}\left\langle
\sum_{1\leq j\neq k,l\neq m\neq n\leq2r}A_{2jk}B_{3lmn}x_{j}x_{k}x_{l}%
x_{m}x_{n}\right\rangle _{3}\\
\left\langle A_{2}B_{3}\right\rangle _{5}  &  =\frac{1}{12}\left\langle
\sum_{1\leq j\neq k,l\neq m\neq n\leq2r}A_{2jk}B_{3lmn}x_{j}x_{k}x_{l}%
x_{m}x_{n}\right\rangle _{5}%
\end{align}%
\begin{align}
\left\langle A_{2}B_{4}\right\rangle _{2}  &  =\frac{1}{48}\left\langle
\sum_{1\leq j\neq k,l\neq m\neq n\neq o\leq2r}A_{2jk}B_{4lmno}x_{j}x_{k}%
x_{l}x_{m}x_{n}x_{o}\right\rangle _{1}\\
\left\langle A_{2}B_{4}\right\rangle _{4}  &  =\frac{1}{48}\left\langle
\sum_{1\leq j\neq k,l\neq m\neq n\neq o\leq2r}A_{2jk}B_{4lmno}x_{j}x_{k}%
x_{l}x_{m}x_{n}x_{o}\right\rangle _{4}\\
\left\langle A_{2}B_{4}\right\rangle _{6}  &  =\frac{1}{48}\left\langle
\sum_{1\leq j\neq k,l\neq m\neq n\neq o\leq2r}A_{2jk}B_{4lmno}x_{j}x_{k}%
x_{l}x_{m}x_{n}x_{o}\right\rangle _{6}%
\end{align}


\subsection{Blades and Versors}

A blade is
\begin{equation}
\mathrm{b}_{p}=v_{1}\wedge\cdots\wedge v_{p}=\left\langle v_{1}\wedge
\cdots\wedge v_{p}\right\rangle _{p}%
\end{equation}
while a versor is
\begin{equation}
\mathrm{v}_{p}=v_{1}\cdots v_{p}=\sum_{0\leq q\leq p}\left\langle v_{1}\cdots
v_{p}\right\rangle _{q},
\end{equation}
hence
\begin{equation}
\mathrm{b}_{p}=\left\langle v_{1}\cdots v_{p}\right\rangle _{p}.
\end{equation}
E.g.%
\begin{equation}
\mathrm{v}_{2}=\sum_{1\leq j\leq k\leq2r}v_{jk}x_{j}x_{k}=\left(  \sum_{1\leq
j\leq2r}v_{jj}\right)  I_{\mathcal{F}}+\sum_{1\leq j<k\leq2r}v_{jk}x_{j}x_{k}.
\end{equation}


\subparagraph{Inverse $^{-1}$}

\subparagraph{%
\begin{equation}
\mathrm{v}_{p}^{-1}=\frac{\mathrm{v}_{p}^{\dagger}}{\left\vert \mathrm{v}%
_{p}\right\vert ^{2}}%
\end{equation}
}

\subparagraph{$\dagger$ is reversion $^{t}$ composed with complex conjugation
\begin{equation}
\mathrm{v}_{p}^{\dagger}=\left(  -1\right)  ^{\frac{p\left(  p-1\right)  }{2}%
}\overline{v_{p}\cdots v_{1}}%
\end{equation}
Grade involution $\ast$}

\subparagraph{%
\begin{equation}
\mathrm{v}_{p}^{\ast}=\left(  -1\right)  ^{p}\mathrm{v}_{p}%
\end{equation}
Scaler product}

\subparagraph{%
\begin{equation}
\left\langle \left.  \mathrm{v}_{p}\right\vert \mathrm{w}_{p}\right\rangle
=\left\langle \mathrm{v}_{p}^{\dagger}\mathrm{w}_{p}\right\rangle _{0}=\bar
{v}_{1}w_{1}\cdots\bar{v}_{p}w_{p}%
\end{equation}
Clifford Valued Norm
\begin{equation}
\left\vert \mathrm{v}_{p}\right\vert ^{2}=\left\langle \left.  \mathrm{v}%
_{p}\right\vert \mathrm{v}_{p}\right\rangle =\left\langle \mathrm{v}%
_{p}^{\dagger}\mathrm{v}_{p}\right\rangle _{0}=\left(  -1\right)
^{\frac{p\left(  p-1\right)  }{2}}\left\vert v_{1}\right\vert ^{2}%
\cdots\left\vert v_{p}\right\vert ^{2}.
\end{equation}
Clifford conjugation}%

\begin{equation}
\mathrm{v}_{p}^{\ddagger}=\mathrm{v}_{p}^{\dagger\ast}=\left(  -1\right)
^{\frac{2p+p\left(  p-1\right)  }{2}}\overline{v_{p}\cdots v_{1}}=\left(
-1\right)  ^{\frac{p\left(  p+1\right)  }{2}}\overline{v_{p}\cdots v_{1}}.
\end{equation}


\subsection{Inner Products}%

\begin{align}
A\rfloor B  &  =\sum_{0\leq p,q\leq2r}\left\langle A_{p}B_{q}\right\rangle
_{p-q}\\
A\lfloor B  &  =\sum_{0\leq p,q\leq2r}\left\langle A_{p}B_{q}\right\rangle
_{q-p}%
\end{align}
In particular
\begin{align}
a\rfloor A_{p}  &  =\frac{1}{2}\left(  aA_{p}-\left(  -1\right)  ^{p}%
A_{p}a\right) \\
a\lfloor A_{p}  &  =\frac{1}{2}\left(  A_{p}a-\left(  -1\right)  ^{p}%
aA_{p}\right)
\end{align}
and
\begin{equation}
\left(  \left.  a\right\vert b\right)  =a\rfloor b=\left\langle
ab\right\rangle _{0}.
\end{equation}


Thin dot%
\begin{equation}
A\cdot B=\sum_{1\leq p,q\leq2r}\left\langle A_{p}B_{q}\right\rangle
_{\left\vert p-q\right\vert }%
\end{equation}
fat dot
\begin{equation}
A\bullet B=\sum_{0\leq p,q\leq2r}\left\langle A_{p}B_{q}\right\rangle
_{\left\vert p-q\right\vert }.
\end{equation}
and Trace%
\begin{equation}
\operatorname*{Trace}\left\{  A\right\}  =\left\langle A\right\rangle _{0}.
\end{equation}


\subsection{Outer Product}%

\begin{equation}
A\wedge B=\sum_{0\leq p,q\leq2r}\left\langle A_{p}B_{q}\right\rangle _{q+p}%
\end{equation}
In particular
\begin{equation}
A_{4}\wedge A_{4}=\left\langle A_{4}A_{4}\right\rangle _{8}%
\end{equation}


\subsection{Scalar Product}%

\begin{align}
\left\langle \left.  A\right\vert B\right\rangle  &  =\left\langle A^{\dagger
}B\right\rangle _{0}=\sum_{0\leq p\leq2r}\left\langle \left.  A_{p}\right\vert
B_{p}\right\rangle =\sum_{0\leq p\leq2r}A_{p}^{\dagger}\rfloor B_{p}%
=\sum_{0\leq p\leq2r}A_{p}^{\dagger}\lfloor B_{p}\\
\left\vert A\right\vert ^{2}  &  =\left\langle \left.  A\right\vert
A\right\rangle \text{ (which could be positive zero or negative)}%
\end{align}


\subsection{Adjoint}

Let linear transformation $F$
\begin{equation}
F:\mathcal{F}_{1}\rightarrow\mathcal{F}_{1}%
\end{equation}
then it can be naturally extended to
\begin{equation}
F:\mathcal{F}\rightarrow\mathcal{F}%
\end{equation}
and is an algebra homomorphism wrt geometric product $\wedge$ i.e.%

\begin{equation}
F\left(  A\wedge B\right)  =F\left(  A\right)  \wedge F\left(  B\right)
\,\forall A,B\in\mathcal{F}%
\end{equation}
and if $F$ is an isometry i.e.
\begin{equation}
F\left(  A\rfloor B\right)  =F\left(  A\right)  \rfloor F\left(  B\right)
\forall A,B\in\mathcal{F}%
\end{equation}
then
\begin{equation}
F\left(  AB\right)  =F\left(  A\right)  F\left(  B\right)  \,\forall
A,B\in\mathcal{F}.
\end{equation}
The converse is also true i.e.
\begin{equation}
F\left(  AB\right)  =F\left(  A\right)  F\left(  B\right)  \,\forall
A,B\in\mathcal{F\Rightarrow}F\left(  A\rfloor B\right)  =F\left(  A\right)
\rfloor F\left(  B\right)  \forall A,B\in\mathcal{F}.
\end{equation}


The adjoint $\mathcal{F}^{\blacksquare}$ is defined by%
\begin{equation}
\left\langle \left.  \mathcal{F}^{\blacksquare}\left(  A\right)  \right\vert
B\right\rangle =\left\langle \left.  A\right\vert F\left(  B\right)
\right\rangle \forall B\in\mathcal{F}.
\end{equation}
Let
\begin{align}
A_{p}  &  =a_{1}\wedge\cdots\wedge a_{p}\\
B_{p}  &  =b_{1}\wedge\cdots\wedge b_{p}%
\end{align}
then
\begin{align}
\left\langle \left.  A_{p}\right\vert B_{p}\right\rangle  &  =\left(
-1\right)  ^{\frac{p\left(  p-1\right)  }{2}}\left\langle \left.  a_{1}%
\wedge\cdots\wedge a_{p}\right\vert b_{1}\wedge\cdots\wedge b_{p}\right\rangle
=\left(  -1\right)  ^{\frac{p\left(  p-1\right)  }{2}}\left\langle \left(
a_{1}\wedge\cdots\wedge a_{p}\right)  ^{\dagger}\left(  b_{1}\wedge
\cdots\wedge b_{p}\right)  \right\rangle _{0}\\
&  =\left(  -1\right)  ^{\frac{p\left(  p-1\right)  }{2}}\sum_{\substack{1\leq
j_{1}<\cdots<j_{p}\leq2r\\1\leq k_{1}<\cdots<k_{p}\leq2r}}\left\langle
a_{pj_{1}}\cdots a_{1j_{1}}\overline{x_{j_{p}}\wedge\cdots\wedge x_{j_{1}}%
}b_{pk_{p}}\cdots b_{1k_{1}}x_{k_{1}}\wedge\cdots\wedge x_{k_{p}}\right\rangle
_{0}\\
&  =\left(  -1\right)  ^{\frac{p\left(  p-1\right)  }{2}}\sum_{\substack{1\leq
j_{1}<\cdots<j_{p}\leq2r\\1\leq k_{1}<\cdots<k_{p}\leq2r}}\left\vert
\begin{array}
[c]{ccc}%
\left\langle \left.  a_{j_{1}}\right\vert b_{k_{1}}\right\rangle  & \cdots &
\left\langle \left.  a_{j_{1}}\right\vert b_{k_{p}}\right\rangle \\
\vdots & \vdots & \vdots\\
\left\langle \left.  a_{j_{p}}\right\vert b_{k_{1}}\right\rangle  & \cdots &
\left\langle \left.  a_{j_{p}}\right\vert b_{k_{p}}\right\rangle
\end{array}
\right\vert I_{\mathcal{F}}%
\end{align}


An element $A$ of a Clifford Algebra is positive i.e. $A\geq0$
$\operatorname*{Iff}$
\begin{equation}
A=B^{\dagger}B.
\end{equation}
Thus
\begin{equation}
A\geq0\Rightarrow A=A^{\dagger}%
\end{equation}
Note that the square of an operator need not be postive i.e. $C^{2}$ could be
negative, zero or complex (think of complex eigenvalues)!

\subsubsection{Outer Nilpotency}

One can note that $\wedge$ is the adjoint of $\rfloor$
\begin{equation}
\left(  A\wedge B\right)  \ast C=A\ast\left(  B\rfloor C\right)  .
\end{equation}
Thus
\begin{equation}
\left(  A\wedge A\right)  \ast C=A\ast\left(  A\rfloor C\right)
\end{equation}
and if
\begin{equation}
\left(  A_{4}\wedge A_{4}\right)  =0
\end{equation}
then
\begin{equation}
\left\langle \left.  A_{4}\right\vert \left(  A_{4}\rfloor C\right)
\right\rangle =\left\langle \left.  A_{4}\right\vert \left\langle \left(
A_{4}\rfloor C\right)  \right\rangle _{4}\right\rangle =0\forall C
\end{equation}
as%
\begin{equation}
A_{4}\rfloor C=\sum_{0\leq q\leq2r}\left\langle A_{4}C_{q}\right\rangle
_{q-4}=\left\langle A_{4}C_{4}\right\rangle _{0}+\left\langle A_{4}%
C_{5}\right\rangle _{1}+\left\langle A_{4}C_{6}\right\rangle _{2}+\left\langle
A_{4}C_{7}\right\rangle _{3}+\left\langle A_{4}C_{8}\right\rangle _{4}%
+\cdots+\left\langle A_{4}C_{2r}\right\rangle _{2r-4}%
\end{equation}
hence
\begin{equation}
\left\langle A_{4}\rfloor C\right\rangle _{4}=\left\langle A_{4}%
C_{8}\right\rangle _{4}%
\end{equation}
and%
\begin{equation}
\left\langle \left.  A_{4}\right\vert \left(  A_{4}\rfloor C\right)
\right\rangle =\left\langle \left.  A_{4}\right\vert \left\langle \left(
A_{4}\rfloor C_{8}\right)  \right\rangle _{4}\right\rangle =0\forall C_{8}.
\end{equation}
This is equivalent to
\begin{equation}
\left\langle \left.  A_{4}\right\vert \left\langle A_{4}\rfloor x_{k_{1}%
}\wedge\cdots\wedge x_{k_{8}}\right\rangle _{4}\right\rangle =0,\,1\leq
k_{1}<\cdots<k_{8}\leq2r.
\end{equation}%
\begin{align}
A_{4}\rfloor x_{k_{1}}\wedge\cdots\wedge x_{k_{8}}  &  =\left\langle
A_{4}x_{k_{1}}\wedge\cdots\wedge x_{k_{8}}\right\rangle _{4}=\sum_{1\leq
j_{1}<\cdots<j_{4}\leq2r}A_{j_{1}j_{2}j_{3}j_{4}}\left\langle x_{j_{1}%
}x_{j_{2}}x_{j_{3}}x_{j_{4}}x_{k_{1}}x_{k_{2}}x_{k_{3}}x_{k_{4}}x_{k_{5}%
}x_{k_{6}}x_{k_{7}}x_{k_{8}}\right\rangle _{4}\\
&  =A_{k_{1}k_{2}k_{3}k_{4}}x_{k_{5}}x_{k_{6}}x_{k_{7}}x_{k_{8}}-A_{k_{1}%
k_{2}k_{3}k_{8}}x_{k_{4}}x_{k_{5}}x_{k_{6}}x_{k_{7}}+\cdots\\
\left\langle \left.  A_{4}\right\vert \left\langle A_{4}\rfloor x_{k_{1}%
}\wedge\cdots\wedge x_{k_{8}}\right\rangle _{4}\right\rangle  &
=A_{k_{5}k_{6}k_{7}k_{8}}A_{k_{1}k_{2}k_{3}k_{4}}-A_{k_{4}k_{5}k_{6}k_{7}%
}A_{k_{1}k_{2}k_{3}k_{8}}+\cdots=0
\end{align}
I.e.
\begin{equation}
A_{5678}A_{1234}-A_{4567}A_{1238}+\cdots=0
\end{equation}
Let%
\[
A_{4}=A_{1234}e_{1234}+A_{1235}e_{1235}+A_{4678}e_{4678}+A_{5678}e_{5678}%
\]%
\begin{align}
A_{4}\wedge A_{4}  &  =\left(  A_{1234}e_{1234}+A_{1235}e_{1235}%
+A_{4678}e_{4678}+A_{5678}e_{5678}\right)  \wedge\left(  A_{1234}%
e_{1234}+A_{1235}e_{1235}+A_{4678}e_{4678}+A_{5678}e_{5678}\right) \\
&  =\left(  A_{1234}A_{5678}-A_{1235}A_{4678}\right)  e_{12345678}%
\end{align}%
\begin{align}
&  1234,1235,1236,1237,1238,\\
&  1245,1246,1247,1248,\\
&  1345,1346,1348,\\
&  1256,1257,1258,\\
&  1267,1268,\\
&  1278
\end{align}%
\begin{equation}
\binom{4}{1}+\binom{4}{2}+\binom{4}{3}=4+6+4=14
\end{equation}


\subsection{Determinant}

The determinat of $F,\det(F)\in$ $\mathcal{F}$ can be defined by
\begin{equation}
\det(F)\mathbf{I}=F\left(  \mathbf{I}\right)
\end{equation}
or equivantley
\begin{equation}
\det(F)=F\left(  \mathbf{I}\right)  ^{\bot}%
\end{equation}


\subsection{Dual (Hodge)}%

\begin{equation}
A^{\bot}=A\rfloor\mathbf{I}^{-1}=A\mathbf{I}^{-1}%
\end{equation}
where
\begin{equation}
\mathbf{I}^{-1}=\frac{\mathbf{I}^{\dagger}}{\left\vert \mathbf{I}\right\vert
^{2}}%
\end{equation}
In particular
\begin{equation}
\left(  v_{1}\wedge\cdots\wedge v_{p}\right)  ^{\bot}=v_{p+1}\wedge
\cdots\wedge v_{2r}%
\end{equation}


\section{States}

\subsection{Definition}

\qquad A state is a positive normalized linear functional
\begin{align}
f  &  :\mathcal{F\rightarrow}\mathbb{R}\\
f\left(  I_{\mathcal{F}}\right)   &  =1.
\end{align}
It is extremal if it cannot be decomposed into a convex sum of two other
states, any state can be expressed as a convex sum of extremal states i.e.
\begin{equation}
f=\sum_{1\leq j\leq p}\lambda_{j}f_{j}.
\end{equation}


\subsection{Density Operators}

A state $f$ is a Normal state if there exists a density operator
$D=D^{\dagger}\in\mathcal{F}$ such that
\begin{equation}
f\left(  A\right)  =\operatorname*{tr}\left\{  A^{\dagger}D\right\}  \,\forall
A\in\mathcal{F}%
\end{equation}%
\begin{equation}
\left\langle \left.  A^{\dagger}\right\vert D\right\rangle =\sum_{0\leq
j\leq2^{4r}}\left\langle A_{j}^{\dagger}D_{j}\right\rangle _{0}=\sum_{0\leq
j\leq2^{2r}}\operatorname*{tr}\left\{  \left\langle A_{j}^{\dagger}%
D_{j}\right\rangle _{0}\right\}  I_{\mathcal{F}}%
\end{equation}
thus
\begin{equation}
f\left(  A\right)  =\sum_{0\leq j\leq2^{2r}}\operatorname*{tr}\left\{
\left\langle A_{j}^{\dagger}D_{j}\right\rangle _{0}\right\}  .
\end{equation}
In particular as
\begin{equation}
D=\sum_{0\leq j\leq2^{2r}}\left\langle D\right\rangle _{j}\equiv\sum_{0\leq
j\leq2^{2r}}D_{j}%
\end{equation}
one has%
\begin{equation}
\left\langle I_{\mathcal{F}}D\right\rangle _{0}=D_{0}%
\end{equation}
and
\begin{equation}
f_{D}\left(  I_{\mathcal{F}}\right)  =D_{0}=1,
\end{equation}
thus
\begin{equation}
D=I_{\mathcal{F}}+\sum_{1\leq j\leq2^{1r}}D_{j}.
\end{equation}


In particular
\begin{align}
\left\vert \varphi\right\rangle \left\langle \varphi\right\vert  &
=I_{\mathcal{F}}+\sum_{1\leq j\leq2^{2r}}\left\langle \left\vert
\varphi\right\rangle \left\langle \varphi\right\vert \right\rangle _{j}\\
\operatorname*{Trace}\left\{  \left\vert \varphi\right\rangle \left\langle
\varphi\right\vert \right\}    & =\operatorname*{Trace}\left\{  I_{\mathcal{F}%
}\right\}  =\operatorname{Tr}\left\{  I_{\mathcal{F}}\right\}  I_{\mathcal{F}%
}\\
\operatorname{Tr}\left\{  \left\vert \varphi\right\rangle \left\langle
\varphi\right\vert \right\}    & =\operatorname{Tr}\left\{  I_{\mathcal{F}%
}\right\}  =2^{2r}\\
\operatorname*{tr}\left\{  \left\vert \varphi\right\rangle \left\langle
\varphi\right\vert \right\}    & =2^{-2r}\operatorname{Tr}\left\{
I_{\mathcal{F}}\right\}  =1.
\end{align}


Projectors can have $0\leq\left\langle \left\langle p\right\rangle
\right\rangle _{0}\leq1,$ while density operators always have $\left\langle
\left\langle D\right\rangle \right\rangle _{0}=1.$ 

\subsection{Expectation Values and States}

The standard Hamiltonian of particle conserved molecular and atomic systems is
of the form
\begin{equation}
H=\eta_{0}I_{\mathcal{F}}+\sum_{1\leq j_{1}<j_{2}\leq2r}\eta_{2j_{1}j_{2}%
}x_{j_{1}}x_{j_{2}}+\sum_{1\leq j_{1}<j_{2}<j_{3}<j_{4}\leq2r}\eta
_{4j_{1}j_{2}j_{3}j_{4}}x_{j_{1}}x_{j_{2}}x_{j_{3}}x_{j_{4}}%
\end{equation}
and energy in a normal state $f_{D}$ is given by
\begin{equation}
f_{D}\left(  H\right)  =\operatorname*{tr}\left\{  \left\langle H_{0}%
D_{0}\right\rangle _{0}\right\}  +\operatorname*{tr}\left\{  \left\langle
H_{2}D_{2}\right\rangle _{0}\right\}  +\operatorname*{tr}\left\{  \left\langle
H_{4}D_{4}\right\rangle _{0}\right\}  =\operatorname*{tr}\left\{  H^{\dagger
}D\right\}  .
\end{equation}
If the state $D$ is pure then
\begin{equation}
D^{2}=D,
\end{equation}
while if it is mixed (ensemble) then
\begin{gather}
D\geq D^{2}\\
\Updownarrow\\
D=\sum_{1\leq j\leq q}\alpha_{j}D_{j},
\end{gather}
where $\left\{  D_{j}\right\}  $ are pure density operators and
\begin{equation}
\sum_{1\leq j\leq q}\alpha_{j}=1;\qquad0\leq\alpha_{j}<1.
\end{equation}


\subsubsection{The Number Operator}

If a state describes on the average $N$ particles then%
\begin{equation}
f_{D}\left(  \mathcal{N}\right)  =\operatorname*{tr}\left\{  \left\langle
\mathcal{N}_{0}D_{0}\right\rangle _{0}\right\}  +\operatorname*{tr}\left\{
\left\langle \mathcal{N}_{2}D_{2}\right\rangle _{0}\right\}  =N,
\end{equation}
where the number operator $\mathcal{N}$ is
\begin{equation}
\mathcal{N}=\mathcal{N}_{0}+\mathcal{N}_{2}=\mathcal{N}=\frac{1}{2}\left(
\operatorname*{Tr}\left\{  \mathcal{N}\right\}  \right)  I_{\mathcal{F}}%
+i\sum_{1\leq j\leq r}x_{j}x_{j+r}.
\end{equation}
If it describes $N$ particles then
\begin{gather}
\left[  \mathcal{N},D\right]  =0\,\&\,\operatorname*{tr}\left\{  \left\langle
\mathcal{N}D\right\rangle _{0}\right\}  =N\nonumber\\
\Updownarrow\nonumber\\
\operatorname*{tr}\left\{  \left\langle \mathcal{N}^{2}D\right\rangle
_{0}\right\}  -\left(  \operatorname*{tr}\left\{  \left\langle \mathcal{N}%
D\right\rangle _{0}\right\}  \right)  ^{2}=0.
\end{gather}


The number operator expressed in the $\left\{  x_{j}\right\}  $ bases
involves
\begin{equation}
\operatorname*{Tr}\left\{  \mathcal{N}\right\}  =1+2+3+\cdots+r,
\end{equation}
which is infinity when $r=\infty,$ however the back transformation works as
one gets $\infty-\infty$ in the form of $r-r$ as can be seen from%
\begin{align}
\mathcal{N} &  =\frac{1}{2}\left(  \operatorname*{Tr}\mathcal{N}\right)
I_{\mathcal{F}}+i\sum_{1\leq j\leq r}\frac{1}{\sqrt{2}}\left(  a_{j}^{\dagger
}+a_{j}\right)  \frac{i}{\sqrt{2}}\left(  a_{j}^{\dagger}-a_{j}\right)
\nonumber\\
&  =\frac{1}{2}\left(  \operatorname*{Tr}\mathcal{N}\right)  I_{\mathcal{F}%
}+i\sum_{1\leq j\leq r}\frac{i}{2}\left(  -a_{j}^{\dagger}a_{j}+a_{j}%
a_{j}^{\dagger}\right)  \nonumber\\
&  =\frac{1}{2}\left(  \operatorname*{Tr}\mathcal{N}\right)  I_{\mathcal{F}%
}-\frac{1}{2}\sum_{1\leq j\leq r}\left(  -a_{j}^{\dagger}a_{j}+I_{\mathcal{F}%
}-a_{j}^{\dagger}a_{j}\right)  \nonumber\\
&  =\frac{1}{2}\left(  \operatorname*{Tr}\mathcal{N}\right)  I_{\mathcal{F}%
}-\sum_{1\leq j\leq r}\left(  \frac{1}{2}I_{\mathcal{F}}-a_{j}^{\dagger}%
a_{j}\right)  \nonumber\\
&  =\frac{1}{2}\left(  \operatorname*{Tr}\mathcal{N}\right)  I_{\mathcal{F}%
}-\frac{1}{2}\left(  \operatorname*{Tr}\mathcal{N}\right)  I_{\mathcal{F}%
}+\sum_{1\leq j\leq r}a_{j}^{\dagger}a_{j}\nonumber\\
&  =\mathcal{N}.
\end{align}


\subsection{Primitive Idempotents}

If an idempotent $p$ is self adjoint
\begin{equation}
p=p^{\dagger}%
\end{equation}
it is thus positive as
\begin{equation}
p=p^{2}=p^{\dagger}p.
\end{equation}
It is possible to have non selfadjoint idempotents (theses are associated with
non orthogonal projectors and nonorthogonal subspace decompositions), which
are non positive.

If $\left\{  p_{j}\right\}  $ is a set of primitive self adjoint idempotents
of $\mathcal{F}$ forming a resolution of the identity then
\begin{equation}
I_{\mathcal{F}}=\sum_{1\leq j\leq p}p_{j}%
\end{equation}
and
\begin{equation}
f\left(  I_{\mathcal{F}}\right)  =1=\sum_{1\leq j\leq p}f\left(  p_{j}\right)
=\sum_{1\leq j,k\leq p}\lambda_{k}f_{k}\left(  p_{j}\right)  \Leftrightarrow
f_{k}\left(  p_{j}\right)  =\delta_{jk}%
\end{equation}
and thus there is a $1-1$ correspondence between extremal states and primitive
idempotents of $\mathcal{F}.$

A primitive idempotent could be of any order i.e.
\begin{equation}
p=\sum_{0\leq j\leq2^{2r}}\left\langle p\right\rangle _{j}=\sum_{0\leq
j\leq2^{2r}}p_{j},
\end{equation}
where $p_{j}=\left\langle p\right\rangle _{j}$ and is not in general
idempotent. The scalar part can be expressed as
\begin{equation}
\left\langle \left.  p\right\vert I_{\mathcal{F}}\right\rangle =p_{0}%
=\operatorname*{tr}\left\{  p^{\dagger}I_{\mathcal{F}}\right\}  I_{\mathcal{F}%
}=\left\langle p\right\rangle _{0}.
\end{equation}
As $p=p^{2}$ one has in particular that
\begin{align}
p_{0} &  =\left\langle p^{2}\right\rangle _{0}=\left\langle p^{\dagger
}p\right\rangle _{0}\geq0=\sum_{0\leq j\leq2^{2r}}\left\langle p_{j}^{\dagger
}p_{j}\right\rangle _{0}=p_{0}^{2}+\sum_{1\leq j\leq2^{2r}}\left\langle
p_{j}^{\dagger}p_{j}\right\rangle _{0}\\
p_{0}-p_{0}^{2} &  =p_{0}\left(  I_{\mathcal{F}}-p_{0}\right)  =\sum_{1\leq
j\leq2^{2r}}\left\langle p_{j}^{\dagger}p_{j}\right\rangle _{0}\geq0,
\end{align}
so
\begin{equation}
0\leq p_{0}=\operatorname*{tr}\left\{  p_{0}\right\}  I_{\mathcal{F}}\leq
I_{\mathcal{F}}\Longleftrightarrow0\leq\operatorname*{tr}\left\{
p_{0}\right\}  \leq1.
\end{equation}


\subsection{Extension Condition}

As $f\left(  I_{\mathcal{F}}\right)  =1$ one has
\begin{equation}
f\left(  I_{\mathcal{F}}\right)  =\sum_{1\leq j\leq q}f\left(  p_{j}\right)
=1,
\end{equation}
thus
\begin{equation}
0\leq f\left(  p_{j}\right)  \leq1;1\leq j\leq q
\end{equation}
and the necessary and sufficient condition for $f$ to be a state is
\begin{equation}
0\leq f\left(  p\right)  \leq1
\end{equation}
\emph{for all generating sets of self adjoint idempotents }$p.$ This can be
extended to the \emph{Necessary and Sufficient} condition for the
\emph{Representability} of a partial or reduced state, $f$, defined on a
subspace $\mathcal{V\subset F}$:%

\begin{equation}
0\leq f\left(  p\right)  \leq1
\end{equation}
\emph{for all }$\mathcal{V-}$\emph{generating sets of self adjoint idempotents
}$p\in\mathcal{V}.$


\end{document}